\section{Consequences and localized amplitudes}\label{sec:consequences}
We derive optimal truncation and pointwise estimates from the approximation theorems, and then treat Schubert's localized example through the reversed-weight class. The estimates for the original norm are stated in Section~\ref{sec:repair}; the companion-norm consequences are recorded here.

\subsection{Truncation and pointwise consequences}\label{subsec:pc-cor}\label{subsec:dd-cor}
The leading bound in Corollary~\ref{pc:thm13} is Theorem~\ref{pc:main} with $K=N=0$.

\begin{proof}[Proof of Corollary~\ref{pc:cor15}]
The same calculation applies to each of the three amplitude norms. The comparison with the truncation window in \cite{S} is given in Appendix~\ref{sec:changes}.

At $t=0$ the error vanishes. Let $t>0$. Theorem~\ref{pc:main} bounds the error by
\[
  C^NN^{2N}(\alpha\hbar)^{-2N}\delta_K^{K+1}K!\,\pcNorm a,
  \qquad \delta_K=\frac{C\hbar t}{\alpha^2}.
\]
On the stated window $K\ge1$. Stirling's inequality and $K\delta_K\le1-\epsilon/2$ give
\[
  \delta_K^{K+1}K!
  \le\ue\,\delta_K K^{1/2}\ue^{-K}(K\delta_K)^K
  \le\ue\,\delta_K K^{1/2}\ue^{-K\epsilon/2}\ue^{-K}.
\]
Since $K^{1/2}\ue^{-K\epsilon/2}\le(\ue\epsilon)^{-1/2}$ and
$K\ge(1-\epsilon)/\delta_K$, this is at most
\[
  \ue^{1/2}\epsilon^{-1/2}\delta_K
  \exp\Bigl(-\frac{(1-\epsilon)\alpha^2}{C\hbar t}\Bigr).
\]
Finally, the time window, $\alpha^3\ge\kappa\hbar$ and the cap on $\delta$ give
\[
  \frac{\hbar t}{\alpha^2}\le\delta\hbar^{1/2}\alpha^{-3/2}
  \le\delta\kappa^{-1/2}\le\frac1{\max(1,C)}.
\]
Thus $\delta_K\le1$ and $C^N\delta_K=C^{N+1}\hbar t/\alpha^2\le C^{N+1}$. Combining the bounds proves the corollary with
$C'=\ue^{1/2}$.
\end{proof}

\begin{proof}[Proof of Corollary~\ref{pc:cor16}]
Let $E$ be the bracket in the statement. By Theorem~\ref{pc:main} at $N$ and at $N+1$, $u:=\Delta^NE$ lies in $\operatorname{Dom}(\Delta_M)$. Apply
Lemma~\ref{pc:g43} to $u$, and use $N^{2N}\le(\ue/2)^{2N}(2N)!$, $(N+1)^{2N+2}\le54^{N+1}(2N)!$ and $\alpha\hbar\le1$. This proves the estimate for $\alpha\hbar\le1$; the completion argument in Section~\ref{subsec:pc-all} supplies the factor depending on $A$ for the remaining range.
For the bounded-geometry specialization, the balls $B(\gamma\tilde z,r_0)$ are disjoint.
Those with $d(\tilde z,\gamma\tilde z)<2$ lie in $B(\tilde z,2+r_0)$, so comparison of hyperbolic areas gives
$m(z)\le C(r_0)$. This proves the bounded-geometry assertion. \qedhere
\end{proof}

\begin{remark}[pointwise bounds and cusp height]\label{pc:remcusp}
In an embedded cusp of height $H$, deep enough, $m(z)\asymp H$, so the pointwise loss grows like $H^{1/2}$. Theorem~\ref{cl:main} shows that this exponent is necessary on the modular surface for a phase attached to the cusp. The proof above uses one Laplacian and gives the exponent $-2N-2$ of $\alpha\hbar$; we state $-2N-4$ to facilitate the comparison in Appendix~\ref{sec:changes}.
\end{remark}

\begin{corollary}[optimal truncation for the companion norms]\label{cor:dd15}\label{cor:dd15-appx}
Under the hypotheses of Theorem~\ref{thm:dd14}, respectively Theorem~\ref{thm:dd-appx}, the estimate and truncation window
of Corollary~\ref{pc:cor15} hold with $\|a\|_{\mathrm S}$, respectively $\|a\|_{\mathrm W}$, in place of $\|a\|_{\mathrm P}$.
Use the constants of the respective theorem and decrease $\delta$ so that $\delta\le\kappa^{1/2}/\max(1,C)$.
\end{corollary}
\begin{proof}
Apply the calculation in the proof of Corollary~\ref{pc:cor15} to the respective approximation theorem.
The same absolute constant $C'=\ue^{1/2}$ works.
\end{proof}

\begin{corollary}[pointwise approximation for the companion norms]\label{cor:dd16}\label{cor:dd16-appx}
Under the hypotheses of Theorem~\ref{thm:dd14}, respectively Theorem~\ref{thm:dd-appx}, assume that $M$ has injectivity
radius at least $r_0>0$. The bounded-geometry estimate of Corollary~\ref{pc:cor16} holds with $\|a\|_{\mathrm S}$,
respectively $\|a\|_{\mathrm W}$, in place of $\|a\|_{\mathrm P}$, and a constant depending also on $r_0$.
Optimal truncation gives the corresponding pointwise exponential estimate from Corollary~\ref{cor:dd15}.
\end{corollary}
\begin{proof}
Apply Lemma~\ref{pc:g43} to $\Delta^NE$ and use the respective approximation theorem at $N$ and $N+1$, as in the proof of
Corollary~\ref{pc:cor16}. Apply the same argument to the optimal-truncation bounds for the exponential estimates.
On an arbitrary quotient the proof gives the additional factor $m(z)^{1/2}$.
\end{proof}


\subsection{Schubert's example amplitudes}\label{subsec:dd-ex}

\begin{proof}[Proof of Remark~\ref{rem:amplitude-examples}]
The spectral calculus gives $\ue^{\alpha\Lambda}\ue^{-\alpha\Lambda}f=f$ for $f\in L^2$ and
$a=\ue^{-\alpha\Lambda}\ue^{\alpha\Lambda}a$ for $a\in\operatorname{Dom}(\ue^{\alpha\Lambda})$.
These identities give the printed-norm assertions, and multiplication by $\ue^{\beta\langle d\rangle}$ gives the
reversed-weight norm identity. For the example, $\langle d\rangle\le d$ and polar coordinates give
\[
  \norm{\ue^{\beta\langle d\rangle}\ue^{-qX_0}}^2
  \le2\pi\int_0^\infty\ue^{2\beta r-2q\cosh r}\sinh r\,dr<\infty.
\]
The integral is finite since $\cosh r\ge\ue^r/2$ and $\sinh r\le\ue^r$.
The smoothed amplitude is nonzero because the multiplier $\ue^{-\alpha\lambda}$ is strictly positive.
\end{proof}

\begin{lemma}[analytic norm of the localized example]\label{lem:ex}
Fix $r_1<\tfrac12$. There are constants $C_H$ and $C_{\mathrm E}$, depending only on $r_1$, such that for $\mu\ge1$ and
$0<\alpha\le\min(r_1,\mu^{-1/2})$,
\[
\norm{\ue^{\alpha\Lambda}\ue^{-\mu\langle d\rangle}}_{L^2(\Hb)}\le(\pi\ue)^{1/2}\,\ue^{C_H\mu\alpha^2}\norm{\ue^{-\mu\langle d\rangle}}_{L^2(\Hb)}\le C_{\mathrm E}\,\mu^{-1/2}.
\]
\end{lemma}

\begin{proof}
We use the notation of Section~\ref{subsec:gutzmer}, with $x_0=i$.

\emph{The function $H$.} Let $G(z):=-\arccos(z)^2$ and $H(z):=(1+G(z))^{1/2}-1$. The function $G$ is holomorphic on
$\mathbb C\setminus(-\infty,-1]$ and equals $d^2$ at $z=\cosh d$, so $\langle d\rangle=H(X_0)$ on $\Hb$. The number $1+G(z)$ is real and at most $0$ only for
real $z\le\cos1$, so $H$ is holomorphic on the complement of $(-\infty,\cos1]$, real on $(\cos1,\infty)$ and nondecreasing there.
For $0<\rho<\tfrac12$ let $S_\rho:=\{\operatorname{Re}z\ge\cos2\rho,\ |\arg z|\le2\rho\}$, and fix $r_1<r_1'<r_2<\tfrac12$ depending only on $r_1$.
Since $\cos2r_2>\cos1$, $H$ is holomorphic near $S_{r_2}$. As $|z|\to\infty$ in $S_{r_2}$ we have $G=L^2$ with
$L=\log(z+(z^2-1)^{1/2})=\log(2z)+O(|z|^{-2})$, hence $H=L-1+O(1/|L|)$; so $M_H:=\sup_{S_{r_2}}|H(z)-\log z|<\infty$. There is $c>0$,
depending only on $r_1',r_2$, such that the disc of radius $c|z|$ about any $z\in S_{r_1'}$ lies in $S_{r_2}$, and Cauchy's estimates for
$H-\log$ on these discs give
\[
|H'(z)|\le C_1/|z|,\qquad|H''(z)|\le C_1/|z|^2\qquad(z\in S_{r_1'}).
\]

\emph{The orbit norms.} By Lemma~\ref{lem:gutzmer}(b), on the orbit $G\exp(i\rho h_0)x_0$, $\rho\le r_1'$, the extension of $X_0$ takes the
values $\tilde X_0=cA+isB$ with $c=\cos2\rho$, $s=\sin2\rho$, $A=X_0(gx_0)\ge1$ and $|B|\le A$; these lie in $S_\rho\subset S_{r_1'}$, and so does
the vertical segment from $cA$ to $\tilde X_0$. So $f:=\ue^{-\mu\langle d\rangle}=\ue^{-\mu H(X_0)}$ extends holomorphically to $\Xi_{r_1'}$. Since $H'(cA)$ is
real, Taylor's formula on that segment gives $\operatorname{Re}H(\tilde X_0)\ge H(cA)-\tfrac12\sup|H''|\,s^2B^2\ge H(cA)-C_1s^2/(2c^2)$, and
$H(A)-H(cA)\le(1-c)A\sup_{[cA,A]}|H'|\le C_1(1-c)/c$. With $1-c\le2\rho^2$, $s^2\le4\rho^2$ and $c\ge\cos1>\tfrac12$,
\[
\operatorname{Re}H(\tilde X_0)\ge\langle d\rangle(gx_0)-C_H\rho^2,\qquad C_H:=12\,C_1 .
\]
Hence $|f(g\exp(i\rho h_0)x_0)|\le\ue^{C_H\mu\rho^2}f(gx_0)$ and $N_\rho(f)\le\ue^{C_H\mu\rho^2}\norm f$ for $\rho\le r_1'$. In particular
$\sup_{\rho<r_1'}N_\rho(f)<\infty$, and Lemma~\ref{lem:gutzmer}(a), with $R=r_1'$ and $r=\alpha\le r_1$, gives the first inequality.

\emph{The $L^2$ norm.} In geodesic polar coordinates $\norm f^2=2\pi\int_0^\infty\ue^{-2\mu\langle\rho\rangle}\sinh\rho\,d\rho$, with
$\langle\rho\rangle=(1+\rho^2)^{1/2}-1$. For $\rho\le1$, $\langle\rho\rangle\ge\rho^2/(1+\sqrt2)$ and $\sinh\rho\le2\rho$; for $\rho\ge1$,
$\langle\rho\rangle\ge\rho-1$ and $\sinh\rho\le\ue^\rho$. So for $\mu\ge1$
\[
\norm f^2\le2\pi\Bigl(\int_0^\infty2\rho\,\ue^{-2\mu\rho^2/(1+\sqrt2)}d\rho+\ue\int_0^\infty\ue^{-(2\mu-1)u}du\Bigr)\le\frac{2\pi(2+\ue)}{\mu},
\]
and the second inequality holds with $C_{\mathrm E}:=(\pi\ue)^{1/2}\ue^{C_H}(2\pi(2+\ue))^{1/2}$, because $\mu\alpha^2\le1$.
\end{proof}

\begin{remark}[Consequence for the example of \texorpdfstring{\cite{S}}{[S]}]
\cite[Remark~1.7, (1.38)]{S} gives the example $a_\hbar:=\hbar^{-\delta/2}\ue^{-\hbar^{-\delta}\langle d\rangle}$, said there to lie in $H_{\alpha,\beta}$ with $\beta=1$ and $\alpha=\hbar^\delta$, and to give long times for $\delta<1/3$. In the reversed-weight $\langle d\rangle$-class, $\norm{a_\hbar}^{\rm sw}_{\alpha,\beta}=\hbar^{-\delta/2}\norm{\ue^{\alpha\Lambda}\ue^{-(\hbar^{-\delta}-\beta)\langle d\rangle}}$. Apply Lemma~\ref{lem:ex} with $\mu=\hbar^{-\delta}-\beta$ and $\alpha_\hbar:=(\hbar^{-\delta}-\beta)^{-1/2}\asymp\hbar^{\delta/2}$:
\begin{itemize}
\item $\norm{a_\hbar}^{\rm sw}_{\alpha_\hbar,\beta}\le C$ uniformly in $\hbar$, for every fixed $\beta$ (not only $\beta=1$);
\item the reversed-weight $\langle d\rangle$-class transfers to the $\tilde\psi$-class by Lemma~\ref{lem:m2w}, at the radius $c_0\alpha_\hbar$ with the fixed ratio $c_0$;
\item so Theorem~\ref{thm:dd-appx} (any $\beta>1/2$, every Fuchsian $\Gamma$) applies when $(c_0\alpha_\hbar)^3\ge\kappa\hbar$, i.e. for $0<\delta<2/3$, and for $\hbar$ small enough that $\mu=\hbar^{-\delta}-\beta\ge1$ and $\alpha_\hbar\le r_1$. The case $\delta=0$ is not covered: $\ue^{-\langle d\rangle}$ has infinite $\ue^{\beta\langle d\rangle}$-weighted $L^2$ norm for every $\beta>1/2$. The time window is $t\le\delta'(\alpha_\hbar/\hbar)^{1/2}\asymp\hbar^{\delta/4-1/2}$.
\end{itemize}
The radius $\alpha=\hbar^\delta$ of \cite{S} gave $\delta<1/3$ and $t\ll\hbar^{-(1-3\delta)/2}$; the holomorphic tube bound gives the larger radius $\hbar^{\delta/2}$. Thus Schubert's localized family is covered by the reversed-weight theorem (Theorem~\ref{thm:dd14}), through its power-weight version (Theorem~\ref{thm:dd-appx}), for $0<\delta<2/3$. This application does not require membership in the original-order class of Theorem~\ref{pc:main}.
\end{remark}
